% GENERATED FILE. Edit canonical lessons, then run npm run generate:books.
% canonical-source-hash: fnv1a64:e586dc2172d256ae
% canonical-lessons: JA-C293-hanabi, JA-C293-uchiageru, JA-C293-yatai, JA-C293-yozora, JA-C293-nigiwau

\chapter{Fireworks over the River}
\label{ch:ja-fireworks-over-the-river}
\hlchaptermodality{293}

\begin{chapteropening}
\textbf{By the end of this chapter:} \emph{I can say where the fireworks were and who came.}

Complete the last lesson of chapter 293: I can say where the fireworks were and who came.
\end{chapteropening}

\section[\texorpdfstring{\ja{はなび} (hanabi)}{はなび (hanabi)}]{\ja{はなび} (hanabi) — fireworks}

\label{lesson:JA-C293-hanabi}



\begin{quote}
Before the new one: say the Japanese for to die, then the Japanese for to support.
\end{quote}

\subsection*{You'll want to know: \ja{はなび}}
\textbf{\ja{はなび}} — \emph{hanabi} — \textquotedblleft{}fireworks\textquotedblright{}.

\begin{quote}
\textbf{\ja{なつは} \ja{かわで} \ja{はなびが} \ja{あります。}} — \emph{natsu wa kawa de hanabi ga arimasu} — \textquotedblleft{}In summer there are fireworks by the river.\textquotedblright{}
\end{quote}

\begin{cousinweb}
\textbf{\ja{はなび}} (\emph{hanabi}) is a native Japanese word (wago), inherited rather than borrowed from Chinese, built from \emph{hana}, \textquotedblleft{}flower\textquotedblright{}, and \emph{bi}, \textquotedblleft{}fire\textquotedblright{} (from \emph{hi}) (JMdict, KANJIDIC2).
\end{cousinweb}

\begin{grammarlens}[title={What you've built}]
One more word for telling what happened.
\end{grammarlens}

\subsection*{Your turn}
\begin{itemize}
\raggedright
  \item \emph{Say it:} \emph{hanabi}
  \item \emph{Say it:} \emph{hanabi}, once more
  \item \emph{Say it:} \emph{natsu wa kawa de hanabi ga arimasu} — In summer there are fireworks by the river.
  \item \emph{Recall:} say the Japanese for to die, then the Japanese for to support, then say \emph{hanabi} again
\end{itemize}

\begin{tcolorbox}[breakable,colback=teal!4,colframe=teal!35!black,title={Before you move on}]
What does \ja{はなび} mean? (\textquotedblleft{}Fireworks\textquotedblright{}.) Say it once more.
\end{tcolorbox}
\section[\texorpdfstring{\ja{うちあげる} (uchiageru)}{うちあげる (uchiageru)}]{\ja{うちあげる} (uchiageru) — to set off (fireworks); to launch}

\label{lesson:JA-C293-uchiageru}



\begin{quote}
Before the new one: say the Japanese for to support, then the Japanese for fireworks.
\end{quote}

\subsection*{You'll want to know: \ja{うちあげる}}
\textbf{\ja{うちあげる}} — \emph{uchiageru} — \textquotedblleft{}to set off (fireworks); to launch\textquotedblright{}.

\begin{quote}
\textbf{\ja{かわの} \ja{うえで} \ja{はなびを} \ja{うちあげます。}} — \emph{kawa no ue de hanabi o uchiagemasu} — \textquotedblleft{}They set off fireworks over the river.\textquotedblright{}
\end{quote}

Polite \textbf{\ja{うちあげます}}, \emph{uchiagemasu}; an ichidan verb. Past (polite): \textbf{\ja{まちが} \ja{はなびを} \ja{うちあげました。}}, \emph{machi ga hanabi o uchiagemashita}.

\begin{cousinweb}
\textbf{\ja{うちあげる}} (\emph{uchiageru}) is a native Japanese word (wago), inherited rather than borrowed from Chinese, built from \emph{uchi}, \textquotedblleft{}striking\textquotedblright{} (from \emph{utsu}), and \emph{ageru}, \textquotedblleft{}to raise\textquotedblright{} (JMdict, KANJIDIC2).
\end{cousinweb}

\begin{grammarlens}[title={What you've built}]
One more word for telling what happened.
\end{grammarlens}

\subsection*{Your turn}
\begin{itemize}
\raggedright
  \item \emph{Say it:} \emph{uchiageru}
  \item \emph{Say it:} \emph{uchiageru}, once more
  \item \emph{Say it:} \emph{kawa no ue de hanabi o uchiagemasu} — They set off fireworks over the river.
  \item \emph{Recall:} say the Japanese for to support, then the Japanese for fireworks, then say \emph{uchiageru} again
\end{itemize}

\begin{tcolorbox}[breakable,colback=teal!4,colframe=teal!35!black,title={Before you move on}]
What does \ja{うちあげる} mean? (\textquotedblleft{}To set off (fireworks); to launch\textquotedblright{}.) Say it once more.
\end{tcolorbox}
\section[\texorpdfstring{\ja{やたい} (yatai)}{やたい (yatai)}]{\ja{やたい} (yatai) — a food stall; a street stall}

\label{lesson:JA-C293-yatai}



\begin{quote}
Before the new one: say the Japanese for fireworks, then the Japanese for to set off (fireworks).
\end{quote}

\subsection*{You'll want to know: \ja{やたい}}
\textbf{\ja{やたい}} — \emph{yatai} — \textquotedblleft{}a food stall; a street stall\textquotedblright{}.

\begin{quote}
\textbf{\ja{えきの} \ja{まえに} \ja{やたいが} \ja{ならんでいます。}} — \emph{eki no mae ni yatai ga narande imasu} — \textquotedblleft{}Food stalls stand in a row in front of the station.\textquotedblright{}
\end{quote}

\begin{cousinweb}
\textbf{\ja{やたい}} (\emph{yatai}) is a mixed word, native and Sino-Japanese: \emph{ya}, \textquotedblleft{}house\textquotedblright{}, is native Japanese, and \emph{tai}, \textquotedblleft{}a stand\textquotedblright{}, is a Chinese character read with its on'yomi, a reading borrowed from Chinese (Middle Chinese \emph{doj}) (JMdict, KANJIDIC2; Middle Chinese in Baxter's transcription).
\end{cousinweb}

\begin{grammarlens}[title={What you've built}]
One more word for telling what happened.
\end{grammarlens}

\subsection*{Your turn}
\begin{itemize}
\raggedright
  \item \emph{Say it:} \emph{yatai}
  \item \emph{Say it:} \emph{yatai}, once more
  \item \emph{Say it:} \emph{eki no mae ni yatai ga narande imasu} — Food stalls stand in a row in front of the station.
  \item \emph{Recall:} say the Japanese for fireworks, then the Japanese for to set off (fireworks), then say \emph{yatai} again
\end{itemize}

\begin{tcolorbox}[breakable,colback=teal!4,colframe=teal!35!black,title={Before you move on}]
What does \ja{やたい} mean? (\textquotedblleft{}A food stall; a street stall\textquotedblright{}.) Say it once more.
\end{tcolorbox}
\section[\texorpdfstring{\ja{よぞら} (yozora)}{よぞら (yozora)}]{\ja{よぞら} (yozora) — the night sky}

\label{lesson:JA-C293-yozora}



\begin{quote}
Before the new one: say the Japanese for to set off (fireworks), then the Japanese for a food stall.
\end{quote}

\subsection*{You'll want to know: \ja{よぞら}}
\textbf{\ja{よぞら}} — \emph{yozora} — \textquotedblleft{}the night sky\textquotedblright{}.

\begin{quote}
\textbf{\ja{よぞらに} \ja{はなびが} \ja{ひろがりました。}} — \emph{yozora ni hanabi ga hirogarimashita} — \textquotedblleft{}Fireworks spread across the night sky.\textquotedblright{}
\end{quote}

\begin{cousinweb}
\textbf{\ja{よぞら}} (\emph{yozora}) is a native Japanese word (wago), inherited rather than borrowed from Chinese, built from \emph{yo}, \textquotedblleft{}night\textquotedblright{}, and \emph{zora}, \textquotedblleft{}sky\textquotedblright{} (from \emph{sora}) (JMdict, KANJIDIC2).
\end{cousinweb}

\begin{grammarlens}[title={What you've built}]
One more word for telling what happened.
\end{grammarlens}

\subsection*{Your turn}
\begin{itemize}
\raggedright
  \item \emph{Say it:} \emph{yozora}
  \item \emph{Say it:} \emph{yozora}, once more
  \item \emph{Say it:} \emph{yozora ni hanabi ga hirogarimashita} — Fireworks spread across the night sky.
  \item \emph{Recall:} say the Japanese for to set off (fireworks), then the Japanese for a food stall, then say \emph{yozora} again
\end{itemize}

\begin{tcolorbox}[breakable,colback=teal!4,colframe=teal!35!black,title={Before you move on}]
What does \ja{よぞら} mean? (\textquotedblleft{}The night sky\textquotedblright{}.) Say it once more.
\end{tcolorbox}
\section[\texorpdfstring{\ja{にぎわう} (nigiwau)}{にぎわう (nigiwau)}]{\ja{にぎわう} (nigiwau) — to bustle; to be lively with people}

\label{lesson:JA-C293-nigiwau}



\begin{quote}
Before the new one: say the Japanese for a food stall, then the Japanese for the night sky.
\end{quote}

\subsection*{You'll want to know: \ja{にぎわう}}
\textbf{\ja{にぎわう}} — \emph{nigiwau} — \textquotedblleft{}to bustle; to be lively with people\textquotedblright{}.

\begin{quote}
\textbf{\ja{まつりの} \ja{とき、まちは} \ja{にぎわいます。}} — \emph{matsuri no toki, machi wa nigiwaimasu} — \textquotedblleft{}At festival time the town is bustling.\textquotedblright{}
\end{quote}

Polite \textbf{\ja{にぎわいます}}, \emph{nigiwaimasu}; a godan verb. Past (polite): \textbf{\ja{まちが} \ja{にぎわいました。}}, \emph{machi ga nigiwaimashita}.

\begin{cousinweb}
\textbf{\ja{にぎわう}} (\emph{nigiwau}) is a native Japanese word (wago), inherited rather than borrowed from Chinese; its kanji was attached later as a kun'yomi (JMdict).
\end{cousinweb}

\begin{grammarlens}[title={What you've built}]
One more word for telling what happened.
\end{grammarlens}

\subsection*{Your turn}
\begin{itemize}
\raggedright
  \item \emph{Say it:} \emph{nigiwau}
  \item \emph{Say it:} \emph{nigiwau}, once more
  \item \emph{Say it:} \emph{matsuri no toki, machi wa nigiwaimasu} — At festival time the town is bustling.
  \item \emph{Recall:} say the Japanese for a food stall, then the Japanese for the night sky, then say \emph{nigiwau} again
\end{itemize}

\begin{tcolorbox}[breakable,colback=teal!4,colframe=teal!35!black,title={Before you move on}]
What does \ja{にぎわう} mean? (\textquotedblleft{}To bustle; to be lively with people\textquotedblright{}.) Say it once more.
\end{tcolorbox}
